% \begin{document}
\section{Decorated Programs}

Doing a complete proof in Hoare Logic can feel overly verbose. The ``core'' of a proof consists
of the preconditions and postconditions surrounding every command; with those, it's usually
possible to infer the complete shape of the proof tree.

\emph{Decorating} programs is a way to informally write down a Hoare logic proof of
correctness. The idea is to insert assertion decorations between every ``line'' of the program.
Using this informal evidence, you can reconstruct the full formal proof.

\subsection{Informal Rules}

We'll set out a few rules to check whether a decorated program represents a valid proof using
\emph{local consistency} checks. The rules, as you'll see, are informal reflections of the
formal inference rules for Hoare logic.

\medskip
\noindent\textbf{Skip.} For \textbf{skip}, the precondition and postcondition should be the
same, like this:
\begin{center}
\begin{tabular}{l}
$\{P\}$ \\
$\mathbf{skip}$ \\
$\{P\}$
\end{tabular}
\end{center}

\noindent\textbf{Sequence.} For sequences, a new assertion $R$ appears between the two
commands. The two halves $\{P\}\ c_1\ \{R\}$ and $\{R\}\ c_2\ \{Q\}$ must be (recursively)
locally consistent:
\begin{center}
\begin{tabular}{l}
$\{P\}$ \\
$c_1;$ \\
$\{R\}$ \\
$c_2$ \\
$\{Q\}$
\end{tabular}
\end{center}

\noindent\textbf{Assignment.} Assignments are locally consistent when the precondition is the
same as the postcondition except that it substitutes the assigned expression in for the
variable:
\begin{center}
\begin{tabular}{l}
$\{P[a/x]\}$ \\
$x := a$ \\
$\{P\}$
\end{tabular}
\end{center}

\noindent\textbf{Conditions.} An \textbf{if} is locally consistent when both branches are
locally consistent after adding the branch condition to each:
\begin{center}
\begin{tabular}{l}
$\{P\}$ \\
$\mathbf{if}\ b\ \mathbf{then}$ \\
\qquad $\{P \wedge b\}$ \\
\qquad $c_1$ \\
\qquad $\{Q\}$ \\
$\mathbf{else}$ \\
\qquad $\{P \wedge \neg b\}$ \\
\qquad $c_2$ \\
\qquad $\{Q\}$ \\
$\{Q\}$
\end{tabular}
\end{center}

\noindent\textbf{Loops.} A \textbf{while} command should be decorated with a loop invariant:
\begin{center}
\begin{tabular}{l}
$\{P\}$ \\
$\mathbf{while}\ b\ \mathbf{do}$ \\
\qquad $\{P \wedge b\}$ \\
\qquad $c$ \\
\qquad $\{P\}$ \\
$\{P \wedge \neg b\}$
\end{tabular}
\end{center}

\noindent\textbf{Implication.} To capture the CONSEQUENCE rule, you can always write a (valid)
implication to connect two commands:
\begin{center}
\begin{tabular}{l}
$\{P\} \Rightarrow$ \\
$\{Q\}$
\end{tabular}
\end{center}

\nt{
  It's possible to give a less redundant definition of these rules (less decorated), but this won't have the property of having a pre and post condition for every command.
}

\subsection{Decorating a Program}

These informal rules tell you how to verify whether a decorated program is correct, but they
don't by themselves tell you how to construct those decorations (and thereby a proof). You can
do this almost mechanically---except for loop invariants, which still require some creativity
to concoct.

Here's an example program:
\[
\mathbf{while}\ (0 < \mathtt{y})\ \mathbf{do}\ (
\ \mathtt{x} := \mathtt{x} + 1;\ \mathtt{y} := \mathtt{y} - 1
\ )
\]

The first step is to decide what we want to prove about this program. So we write a
precondition and postcondition for the whole program. Intuitively, the program adds
$\mathtt{y}$ onto the initial value of the variable $\mathtt{x}$, so we'll assert that:

\begin{center}
\begin{tabular}{l}
$\{\mathtt{x} = m \wedge \mathtt{y} = n \wedge 0 \leq n\}$ \\
$\mathbf{while}\ (0 < \mathtt{y})\ \mathbf{do}\ ($ \\
\qquad $\mathtt{x} := \mathtt{x} + 1;$ \\
\qquad $\mathtt{y} := \mathtt{y} - 1$ \\
$)$ \\
$\{\mathtt{x} = m + n\}$
\end{tabular}
\end{center}

This program is a while command, so the next step is to come up with a loop invariant. We'll
set up the structure first. We need an assertion $I$ where we can ``connect'' the loop's
precondition and postcondition to our overall precondition and postcondition using
implications, like this:

\begin{center}
\begin{tabular}{l}
$\{\mathtt{x} = m \wedge \mathtt{y} = n \wedge 0 \leq n\} \Rightarrow$ \\
$\{I\}$ \\
$\mathbf{while}\ (0 < \mathtt{y})\ \mathbf{do}\ ($ \\
\qquad $\{I \wedge 0 < y\}$ \\
\qquad $\mathtt{x} := \mathtt{x} + 1;$ \\
\qquad $\mathtt{y} := \mathtt{y} - 1$ \\
\qquad $\{I\}$ \\
$)$ \\
$\{I \wedge 0 \not< \mathtt{y}\} \Rightarrow$ \\
$\{\mathtt{x} = m + n\}$
\end{tabular}
\end{center}

On every iteration of the loop, the variable $\mathtt{x}$ is the sum we want, $m + n$, less the
current value of $\mathtt{y}$, which is the number of iterations remaining. So we'll define the
invariant $I$ like this:
\[
I ::= (\mathtt{x} = m + n - \mathtt{y}) \wedge 0 \leq \mathtt{y}
\]
The top implication follows because $n - \mathtt{y} = 0$, and the bottom implication is valid
because $0 \not< \mathtt{y}$ and $0 \leq \mathtt{y}$ together imply $\mathtt{y} = 0$.

To finish decorating the program, we need an assertion between the two lines in the body of the
loop. By the rule for assignments, we know that this must be $I[\mathtt{y} - 1/\mathtt{y}]$, or:
\[
P_1 ::= (\mathtt{x} = m + n - (\mathtt{y} - 1)) \wedge (0 \leq \mathtt{y} - 1)
\]
To make the assignment to $\mathtt{x}$ locally consistent, then, its precondition must be
$P_1[\mathtt{x} + 1/\mathtt{x}]$, or:
\[
P_2 ::= (\mathtt{x} + 1 = m + n - (\mathtt{y} - 1)) \wedge (0 \leq \mathtt{y} - 1)
\]
It's straightforward to see that the precondition at the top of the loop, $I \wedge 0 <
\mathtt{y}$, implies this new assertion $P_2$. Now we can write our complete decorated program
using these definitions:

\begin{center}
\begin{tabular}{l}
$\{\mathtt{x} = m \wedge \mathtt{y} = n \wedge 0 \leq n\} \Rightarrow$ \\
$\{I\}$ \\
$\mathbf{while}\ (0 < \mathtt{y})\ \mathbf{do}\ ($ \\
\qquad $\{I \wedge 0 < y\} \Rightarrow$ \\
\qquad $\{P_2\}$ \\
\qquad $\mathtt{x} := \mathtt{x} + 1;$ \\
\qquad $\{P_1\}$ \\
\qquad $\mathtt{y} := \mathtt{y} - 1$ \\
\qquad $\{I\}$ \\
$)$ \\
$\{I \wedge 0 \not< \mathtt{y}\} \Rightarrow$ \\
$\{\mathtt{x} = m + n\}$
\end{tabular}
\end{center}

\section{Weakest Preconditions}

Previously, we discussed the issue of completeness---i.e., whether it is possible to derive
every valid partial correctness specification using the axioms and rules of Hoare logic.
Unfortunately, if treated as a pure deductive system, Hoare logic cannot be complete. To see
why, consider the following partial correctness specifications:
\[
\{\mathbf{true}\}\ \mathbf{skip}\ \{P\}
\qquad
\{\mathbf{true}\}\ c\ \{\mathbf{false}\}
\]
The first is valid if and only if the assertion $P$ is valid while the second is valid if and
only if the command $c$ does not halt.

It turns out that the culprit is the CONSEQUENCE rule,
\[
\frac{\models (P \Rightarrow P') \qquad \{P'\}\ c\ \{Q'\} \qquad \models (Q' \Rightarrow
Q)}{\{P\}\ c\ \{Q\}} \; \text{CONSEQUENCE}
\]
which includes two premises about the validity of implications between the assertions
involved. Proving those implications for our two examples would require proving arbitrary
logical predicates $P$ and proving that arbitrary IMP programs $c$ terminate. By G\"odel's
incompleteness theorem, we know that there is no consistent mathematical proof system that can
do the former. Therefore, Hoare logic cannot be complete in that sense.

On the other hand, Hoare logic does enjoy this property:

\thm{Relative Completeness of Hoare Logic}{
$\forall P, Q \in \mathbf{Assn}, c \in \mathbf{Com}.\ \models \{P\}\ c\ \{Q\}$ implies $\vdash
\{P\}\ c\ \{Q\}$.
}

This result states that Hoare logic is no more incomplete than the language of assertions---i.e., if we
had an oracle that could decide the validity of assertions, then we could decide the validity
of partial correctness specifications. In other words, for any valid Hoare triple, there exists
a proof in Hoare logic---as long as you have some external way to decide the validity of the
premises to the CONSEQUENCE rule.\\

Cook's proof of relative completeness depends on the notion of \emph{weakest liberal
preconditions}. Given a command $c$ and a postcondition $Q$, the weakest liberal precondition
is the weakest assertion $P$ such that $\{P\}\ c\ \{Q\}$ is valid. Here, ``weakest'' means that
any other valid precondition implies $P$. That is, $P$ most accurately describes input states
for which $c$ either does not terminate or ends up in a state satisfying $Q$.

\dfn{Weakest Liberal Precondition}{
Formally, an assertion $P$ is a weakest liberal precondition of $c$ and $Q$ if:
\[
\forall \sigma, I.\ \sigma \models_I P \iff (\mathcal{C}[\![c]\!]\, \sigma) \text{ is
undefined} \ \vee\ (\mathcal{C}[\![c]\!]\sigma) \models_I Q
\]
}

\nt{
  The interpretation is the same to check $ P $ and $ Q  $, meaning the program can't change the set of logical variables.
}

We write $\mathrm{wlp}(c, Q)$ for the weakest liberal precondition of command $c$ and
postcondition $Q$. From left-to-right, the formula above states that $\mathrm{wlp}(c, Q)$ is a
valid precondition: $\models \{P\}\ c\ \{Q\}$. The right-to-left implication says it is the
weakest valid precondition: if another assertion $R$ satisfies $\models \{R\}\ c\ \{Q\}$, then
$R$ implies $P$. It can be shown that weakest liberal preconditions are unique modulo
equivalence.

We can calculate the weakest liberal precondition of a command as follows:
\begin{align*}
\mathrm{wlp}(\mathbf{skip}, P) &= P \\
\mathrm{wlp}(x := a, P) &= P[a/x] \\
\mathrm{wlp}((c_1; c_2), P) &= \mathrm{wlp}(c_1, \mathrm{wlp}(c_2, P)) \\
\mathrm{wlp}(\mathbf{if}\ b\ \mathbf{then}\ c_1\ \mathbf{else}\ c_2, P) &= (b \implies
\mathrm{wlp}(c_1, P)) \wedge (\neg b \implies \mathrm{wlp}(c_2, P))
\end{align*}

The definition of $\mathrm{wlp}(\mathbf{while}\ b\ \mathbf{do}\ c, P)$ is slightly more
complicated---it encodes the weakest liberal precondition for each iteration of the loop. To
give the intuition, first define the weakest liberal precondition for a loop that terminates in
$i$ steps as follows:
\begin{align*}
F_0(P) &= \mathbf{true} \\
F_{i+1}(P) &= (\neg b \implies P) \wedge (b \implies \mathrm{wlp}(c, F_i(P)))
\end{align*}
We can then express the weakest liberal precondition using an infinitary conjunction:
\[
\mathrm{wlp}(\mathbf{while}\ b\ \mathbf{do}\ c, P) = \bigwedge_i F_i(P)
\]
See Winskel Chapter 7 for the details of how to encode the weakest liberal precondition for a
while loop as an ordinary assertion.

To check that our definition is correct, we can prove (how?) that it yields a valid partial
correctness specification:

\mlenma{Correctness of Weakest Preconditions}{
\[
\forall c \in \mathbf{Com}, Q \in \mathbf{Assn}.\ \models \{\mathrm{wlp}(c, Q)\}\ c\ \{Q\}
\text{ and } \forall R \in \mathbf{Assn}.\ \models \{R\}\ c\ \{Q\} \text{ implies } (R \implies
\mathrm{wlp}(c, Q))
\]
}

It is not hard to prove that it also yields a provable specification:

\mlenma{Provability of Weakest Preconditions}{
$\forall c \in \mathbf{Com}, Q \in \mathbf{Assn}.\ \vdash \{\mathrm{wlp}(c, Q)\}\ c\ \{Q\}$
}

Relative completeness follows by a simple argument:

\pf{Sketch}{
Let $c$ be a command and let $P$ and $Q$ be assertions such that the partial correctness
specification $\{P\}\ c\ \{Q\}$ is valid. By Lemma~1 we have $\models P \implies
\mathrm{wlp}(c, Q)$. By Lemma~2 we have $\vdash \{\mathrm{wlp}(c, Q)\}\ c\ \{Q\}$. We conclude
$\vdash \{P\}\ c\ \{Q\}$ using the CONSEQUENCE rule.
}

% \end{document}
